\section{Experimental Methodology}

\subsection{Platform and software}

\begin{table}[t]
  \centering
  \caption{Primary experimental platform.}
  \label{tab:platform}
  \begin{tabularx}{\textwidth}{@{}L{0.27\textwidth}Y@{}}
    \toprule
    Component & Configuration \\
    \midrule
    GPU & NVIDIA RTX PRO 6000 Blackwell Server Edition \\
    Target architecture & \code{sm\_120a}; driver 590.48.01 \\
    Build toolkit & CUDA 13.1 \\
    Profiled runtime & CUDA runtime 13.2 \\
    Baseline software & cuML 26.06 RBF \code{SVC.fit()} \\
    Profiling & Nsight Systems 2025.5.2 and Nsight Compute 2025.4.1 \\
    \bottomrule
  \end{tabularx}
\end{table}

Headline timings are profiler-free \code{SVC.fit()} measurements. Nsys is
used for stage attribution and NCU for selected-kernel mechanism analysis.
Final comparisons use fresh processes and balanced execution order when
available.

\subsection{Datasets and preprocessing}

\begin{table}[t]
  \centering
  \caption{Principal workloads and their role.}
  \label{tab:workloads}
  \begin{tabularx}{\textwidth}{@{}L{0.20\textwidth}L{0.25\textwidth}Y@{}}
    \toprule
    Workload & Shape & Purpose \\
    \midrule
    EMBER probe & $65{,}536\times2{,}381$ & Repeated high-dimensional timing,
      profiling, and local-solver capture \\
    EMBER full & $600{,}000\times2{,}381$ train; 200,000 test & Full-scale
      cache and low-precision gates \\
    SUSY & 4.5M train, $d=18$ & Low-dimensional scaling, row-resident reuse,
      and full KKT certification \\
    HIGGS & 10.5M train, $d=28$ & Low-dimensional full-scale generalization \\
    COD-RNA and Skin & Public low-dimensional datasets & External coverage and
      KKT certificates \\
    \bottomrule
  \end{tabularx}
\end{table}

The EMBER-65K probe is standardized with statistics computed from the probe.
The 200K held-out evaluation uses the same probe statistics because it
evaluates models trained on that probe. Full-official EMBER experiments use
statistics computed from the 600K training set. This difference is intentional
but limits direct comparison between probe and full-official campaigns.

The primary coherent-FP16 configuration uses $C=1$, $\gamma=1/2381$,
\code{tol=1e-3}, a 1,024-variable working set, and a cache capacity of 32,768
logical columns.

\subsection{Baseline and fairness dimensions}

The central performance comparison is between coherent FP16 and a
same-framework direct FP32-cache control. Both retain 32,768 logical columns;
the FP32 cache therefore consumes approximately twice the bytes of the FP16
cache. This design isolates representation and execution effects at equal
reuse capacity, not equal memory budget.

The public configuration for this comparison is
\path{artifact/configs/ember65k-ab.toml}. It names the methods
\code{direct\_fp32} and \code{coherent\_fp16}, the balanced schedule, and the
logical cache capacity without naming a machine or a particular measurement
run.

The held-out quality comparison is different: it compares the saved stock
FP32 model with the coherent-FP16 model under the same independent FP32
scorer. The performance and quality tables therefore do not share an identical
comparator. Its public experiment contract is
\path{artifact/configs/ember65k-quality.toml}, which names
\code{stock\_fp32} and \code{coherent\_fp16} and a common held-out split.

\subsection{Numerical contracts}

\paragraph{Exact-FP32 contract.}
The original FP32 RBF kernel defines the objective. Validation checks global
KKT tolerance, objective, equality feasibility, and box feasibility.

\paragraph{Coherent-low-precision contract.}
The local $K[W,W]$, global $K[X,W]$, and persistent cache consume a consistent
low-precision representation. The solver must reach its internal stopping
criterion and retain held-out predictive quality. It is not required to
reproduce every FP32 dual variable, but it must not be called an exact solution
to the original FP32 objective.

\subsection{Measurement and reporting contract}

End-to-end performance uses profiler-free fresh processes, balanced execution
order, and the complete set of repetitions rather than a selected best run.
Profiler traces are used only to attribute mechanisms and are not substituted
for end-to-end timing. Predictive quality is evaluated with a common
independent scorer, while numerical validation reports the stopping criterion
of the declared precision contract separately from an FP32-kernel audit.

Every public run emits a normalized artifact bundle. The manifest records the
configuration and source identities without machine-specific paths; the
measurement records retain method, shape, device class, solver parameters,
timing, and stopping diagnostics; and the summary contains paired statistics
and quality comparisons when requested. These machine-readable files, rather
than an opaque run name or prose log, are the unit used to compare runs.
